\chapter{Comparison-Preserving Changes and Information}\label{chap:7}

\section*{Overview of the Chapter}

The question of this chapter is \textbf{by what information the changes preserving a
comparison can be distinguished, and how they can be constructed}.
Chapter~\ref{chap:6} used normalization to gather together the differences between
representations, and obtained a way of comparing the remaining images.
Now we consider the situation of a refactoring, which changes the internal design while
keeping the correspondence of processing and data.

For example, we extract the update logic from the processing that edits an order into a
shared component, and split the shipping address and the payment information contained in a
single order record into separate data types.
That the shipping address can still be displayed correctly after the change does not tell us
that the payment information is not mixed up when the shipping address is updated.
What is needed is that moving the old data into the new format and then updating, and updating
in the old format and then moving, give the same state of the order.
When only the display on the screen and the call relations are compared, this carrying over
of data can drop out of the check.

In this chapter we single out the changes under which chosen states and structures can be put
into one-to-one correspondence, and study the conditions under which this correspondence
preserves reads, updates, and the carrying over by the processing.

There are two tasks here. One is to \textbf{determine} whether a given change preserves the
comparison.
The other is to \textbf{construct}, from a change of the summarized design, a change of the
detailed design that preserves the comparison.
In this chapter we exhibit, for the normalization of AAT, cases where compatible changes can
be constructed and yet determining a given change requires the information lost in the
normalization.
On that basis, we classify all the options for compatible changes by a group action.

\begin{xltabular}{\linewidth}{@{}LLL@{}}
\toprule
Question & Construction & Outcome \\
\midrule
\endhead
What are the changes that preserve a comparison? & The comparison-preserving group and the
projections to the endpoints & Existence conditions and freedom for a change following a
change on one side (\S\S\ref{sec:7.1}--\ref{sec:7.2}) \\
\addlinespace[3pt]
Can compatibility be determined by observations alone? & The kernel of the observation
homomorphism & Conditions on the information necessary and sufficient for the determination,
and classification of the changes with the same observation (\S\ref{sec:7.3}) \\
\addlinespace[3pt]
Can a change after normalization be lifted to the original? & Restriction to the idempotent
image and a section of it & General conditions for the determination, and construction of
lifts for the generated comparisons of AAT (\S\S\ref{sec:7.4}--\ref{sec:7.6}) \\
\addlinespace[3pt]
How do the operations of software constrain changes? & Connected components of the graph of
operations & Classification and number of the changes that preserve operations
(\S\S\ref{sec:7.7}--\ref{sec:7.8}) \\
\bottomrule
\end{xltabular}

\section{Changes preserving a single comparison}\label{sec:7.1}

We first confirm that, even when the call relations of the processing are the same, a
refactoring can break the data that ought to be carried over.
We then gather the condition of preserving the correspondence with operations and adapters
into a group.

\begin{example}[Extraction into a shared component changes what is carried over]\label{ex:7.1}

We extract the update of the shipping address of an order into a shared component.
Suppose there is a condition that this update does not change the payment category.
Focusing on the internal code $\{0,1\}$ of the payment category, we gather the processing
into a small state machine.
We take a single control point, and let the operation that returns to the same control point
before and after the update be a named self-loop $a$.
We define the correct implementation $X$ and the implementation $Y$ whose conversion in the
extracted component inverts this code by
\begin{equation*}
T_a^X(k)=k,\qquad T_a^Y(k)=1-k
\tag{7.1}\label{eq:7.1}
\end{equation*}
If we read only the vertices and the edges, both have the same update operation.
Updating an order with code 0 once, however, leaves it 0 in $X$ and turns it into 1 in $Y$.
Reading the Law ``an update of the shipping address keeps the payment category'' lets us
distinguish this difference.
In addition to the agreement of the call relations, we need to read the conditions on the
data that the processing carries over.

Likewise for changes, the part remaining in the observation can agree while the
compatibility with operations and adapters differs.
To describe this difference, we fix a single comparison.

\end{example}

\begin{definition}[Comparison-preserving group]\label{def:7.2}

Take a morphism $c:X\to Y$ of a category $E$ and subgroups
$G_X\le\mathrm{Aut}_E(X)$, $G_Y\le\mathrm{Aut}_E(Y)$ of the permitted endpoint changes.
We take the group of changes preserving $c$ to be
\begin{equation*}
\Gamma_c=\{(u,v)\in G_X\times G_Y\mid vc=cu\}
\tag{7.2}\label{eq:7.2}
\end{equation*}
A change satisfying this equation is said to be \textbf{compatible} with the comparison.
The product of the group is composition of morphisms, and $u_2u_1$ applies $u_2$ after
$u_1$.

In a refactoring as well, once a single compatible invertible reference correspondence
linking the old and the new states is fixed, the difference from another invertible
correspondence of states can be expressed as an automorphism of the original set of states.
In \S\ref{sec:7.8} we use this reduction to move the conditions for preserving reads and
operations to conditions on automorphisms.

In AAT we take $E=E_{\mathrm{geom}}$ and write the composite of the two-level projections of
Chapter~\ref{chap:1} as $\pi:E_{\mathrm{geom}}\to B$.
Besides the case where all automorphisms are permitted, we use the changes fixing the base
\begin{equation*}
G_X^0=\ker\bigl(\mathrm{Aut}(X)\xrightarrow{\pi}
\mathrm{Aut}(\pi X)\bigr),\qquad
\Gamma_c^0=\Gamma_c\cap(G_X^0\times G_Y^0)
\tag{7.3}\label{eq:7.3}
\end{equation*}
The unadorned $\Gamma_c$ is the notation for the case where Definition~\ref{def:7.2} is used
with the full endpoint automorphism groups.
What is fixed is the image in the base of the changes $u,v$, and $\pi(c)$ may remain a
general comparison.
Moreover, the condition of fixing the base and the condition of fixing the coefficients can
be specified separately.

\end{definition}

\begin{proposition}[The group obtained from a comparison object of full geometries]\label{prop:7.3}

For a comparison $c:X\to Y$ of full geometries, the group $\Gamma_c^0$ of \eqref{eq:7.3} is
a subgroup of $G_X^0\times G_Y^0$.
Using the embedding $J:E_{\mathrm{geom}}\to\mathrm{Kar}(E_{\mathrm{geom}})$ by the identity
idempotents, it is isomorphic to the group of invertible changes of $J(c)$ fixing the base
in the category of comparisons $\mathcal{M}(E_{\mathrm{geom}})$ of Chapter~\ref{chap:6}.
This identification commutes with the projections to the endpoints.

\end{proposition}

\begin{proof}

The identity pair is compatible, and the product and the inverse of two compatible pairs
also satisfy $vc=cu$.
These operations preserve the identity map of the base as well.
Since the morphisms between the $J(X)=(X,1_X)$ carry no additional constraint from
idempotents, the same endpoint pair can be read as an invertible square in the category of
morphisms of the idempotent completion.
This correspondence and its inverse preserve the morphisms at the endpoints, hence preserve
the group law, the projections, and the condition of fixing the base.

\end{proof}

When only one side is changed, the other need not be able to follow.
When $(u,v)$ preserves the comparison, we say that $v$ follows $u$.
When following is possible, all the candidates can be described by the changes that leave
the comparison as it is.

\begin{theorem}[Projections to the endpoints and following changes]\label{thm:7.4}

We use the same comparison of full geometries and base-fixing groups as in
Proposition~\ref{prop:7.3}.
Take the projections to the endpoints $p_X:\Gamma_c^0\to G_X^0$ and
$p_Y:\Gamma_c^0\to G_Y^0$.
Their kernels are, respectively,
\begin{equation*}
\ker p_X\cong T_c:=\{v\in G_Y^0\mid vc=c\},\qquad
\ker p_Y\cong S_c:=\{u\in G_X^0\mid cu=c\}
\tag{7.4}\label{eq:7.4}
\end{equation*}
A necessary and sufficient condition for the existence of a $v$ following $u\in G_X^0$ is
$u\in\mathrm{im}\,p_X$.
If one following change $v_0$ exists, all the candidates are $T_cv_0$.
Likewise, if a $u_0$ following $v\in G_Y^0$ exists, all the candidates are $S_cu_0$.

\end{theorem}

\begin{proof}

The first kernel is obtained by substituting $u=1_X$ into \eqref{eq:7.2}. The second is
similar.
From $v_0c=cu$ we get $v_0^{-1}c=cu^{-1}$. Using this together with $vc=cu$, we obtain
\begin{equation*}
(vv_0^{-1})c=vc\,u^{-1}=c
\tag{7.5}\label{eq:7.5}
\end{equation*}
Hence we may write it uniquely as $v=tv_0$, $t=vv_0^{-1}\in T_c$.
Conversely, if $t\in T_c$ then $tv_0c=tcu=cu$.
For the other side the computation is the same, using $cu_0^{-1}=v^{-1}c$.

\end{proof}

When the action of a group on a nonempty set is free and transitive, the set is called a
\textbf{torsor} under the group.
That is, between two candidates there is exactly one element of the group carrying one to
the other.
The candidate sets of Theorem~\ref{thm:7.4} are torsors under $T_c$ and $S_c$ respectively,
by composition on the left.
Choosing one candidate as a reference identifies the set with the group, but the candidate
set itself needs no specified reference.

\begin{corollary}[Changes of an isomorphism comparison]\label{cor:7.5}

If a comparison $c:X\xrightarrow{\sim}Y$ of full geometries is an isomorphism, then
\begin{equation*}
\Gamma_c^0=\{(u,cuc^{-1})\mid u\in G_X^0\},\qquad
\Gamma_c^0\cong G_X^0
\tag{7.6}\label{eq:7.6}
\end{equation*}
When the full endpoint automorphism groups are used in an arbitrary category as well, the
comparison-preserving group is isomorphic as a group to $\mathrm{Aut}(X)$ via
$u\mapsto(u,cuc^{-1})$.

\end{corollary}

\begin{proof}

The equation $vc=cu$ is equivalent to $v=cuc^{-1}$.
In the case of the full endpoint groups, this gives both inverses and the group law.
In the base-fixing version for full geometries,
$\pi(cuc^{-1})=\pi(c)1\pi(c)^{-1}=1$ shows that the same correspondence preserves the
base-fixing groups.

\end{proof}

\section{Carrying changes from the same generator}\label{sec:7.2}

In Chapter~\ref{chap:5} we constructed two routes and a comparison from the same generator.
Sharing this generator also answers which changes are to be combined at the two endpoints.

\begin{theorem}[Compatibility of generated changes]\label{thm:7.6}

Take the common original object $S$ of Chapter~\ref{chap:5}, \S\ref{sec:5.8}, the ends
$X,Y$ of the two routes, and the generated comparison $c:X\to Y$.
Let $G_S^0$ be the group of automorphisms of the original object $S$ fixing the base.
Writing the homomorphisms induced by the transports and the pullbacks along the two routes
as $F_*:G_S^0\to G_X^0$ and $G_*:G_S^0\to G_Y^0$, we have
\begin{equation*}
G_*(a)c=cF_*(a)\qquad(a\in G_S^0).
\tag{7.7}\label{eq:7.7}
\end{equation*}
Setting further $J=G_*^{-1}(T_c)$, for independently chosen $a,b\in G_S^0$,
\begin{equation*}
\bigl(F_*(a),G_*(b)\bigr)\in\Gamma_c^0
\quad\Longleftrightarrow\quad ba^{-1}\in J.
\tag{7.8}\label{eq:7.8}
\end{equation*}
Hence, once the change $a$ on one side is fixed, the changes of the source usable on the
other side are $Ja$.

\end{theorem}

\begin{proof}

Write the cartesian legs of the two routes as $q_X:X\to S$ and $q_Y:Y\to S$.
Here, as in Chapter~\ref{chap:5}, we compose the legs after moving to the refinement side
when necessary.
The factorizations at the time of generation and the triangle of the comparison are
$q_XF_*(a)=aq_X$, $q_YG_*(a)=aq_Y$, and $q_Yc=q_X$.
Composing both sides of \eqref{eq:7.7} with $q_Y$ therefore gives $aq_X$ in either case.
Both sides lie over the same morphism of the base. Using the uniqueness coming from
cartesianness of the core and then from cartesianness of the geometry, we obtain
\eqref{eq:7.7}.

By \eqref{eq:7.7} and Theorem~\ref{thm:7.4}, compatibility of $G_*(b)$ with $F_*(a)$ is
equivalent to $G_*(b)G_*(a)^{-1}\in T_c$.
Since these are homomorphisms, this becomes \eqref{eq:7.8}.

\end{proof}

Carrying the same change along both routes gives compatibility, and when different changes
are chosen the remaining difference must lie in $J$.
This description is preserved when the presentation of the generator is changed.

\begin{proposition}[Change of presentation of the input and regeneration]\label{prop:7.7}

Write the diagram consisting of the common original objects of Chapter~\ref{chap:5} as
$(S_v,L_e,\kappa_f)$.
Here $L_e$ is the morphism of an edge and $\kappa_f$ is the comparison specified on a face.
Give at each vertex an isomorphism $\theta_v:S'_v\xrightarrow{\sim}S_v$ fixing the base and
the coefficients, and form the new input by
\begin{equation*}
L'_e=\theta_{t(e)}^{-1}L_e\theta_{s(e)},\qquad
\kappa'_f=\theta_{t(f)}^{-1}\kappa_f\theta_{t(f)}
\tag{7.9}\label{eq:7.9}
\end{equation*}
The choices of transport of Chapter~\ref{chap:4} are also carried through these
isomorphisms, and the two routes are constructed from the new input.
Then, with the endpoint isomorphisms
$b_v:X'_v\xrightarrow{\sim}X_v$ and $d_v:Y'_v\xrightarrow{\sim}Y_v$ given by universality,
\begin{equation*}
\begin{aligned}
c'_v&=d_v^{-1}c_vb_v,\\
F'_*(\theta_v^{-1}a\theta_v)&=b_v^{-1}F_*(a)b_v,\\
G'_*(\theta_v^{-1}a\theta_v)&=d_v^{-1}G_*(a)d_v
\end{aligned}
\tag{7.10}\label{eq:7.10}
\end{equation*}
The comparison-preserving group, the residual subgroup of \eqref{eq:7.8}, and the set of
following changes are carried along this correspondence.

\end{proposition}

\begin{proof}

The edges carried through the isomorphisms are again strongly opcartesian morphisms.
Indeed, composing $\theta$ and its inverse before and after the original factorization
carries over existence and uniqueness.
The same holds for the cartesian morphisms of the pullback.
Hence the universality of each route, old and new, yields $b_v,d_v$ uniquely.
Abbreviating the two legs as $q_X,q_Y$, these satisfy
$q'_X=\theta^{-1}q_Xb$ and $q'_Y=\theta^{-1}q_Yd$.
The candidate $d^{-1}cb$ satisfies the triangle of the new comparison, hence coincides with
$c'$ by uniqueness.
The formulas for the two changes are obtained likewise, by substituting into the new
factorizations.
The conditions on edges and faces are preserved because the intermediate inverses in
\eqref{eq:7.9} cancel.
Finally, substituting \eqref{eq:7.10} into \eqref{eq:7.2} and \eqref{eq:7.8} yields the
correspondence of the groups and of the candidates.

\end{proof}

\section{Conditions for distinguishing by observation}\label{sec:7.3}

To distinguish the changes compatible with a single comparison, it is not always necessary
to read all the components of a change.
Which components may be forgotten can be stated precisely by studying the changes that look
like the identity under the observation.

\begin{definition}[Observation of a change]\label{def:7.8}

Take a group $Q$ of changes, a compatible subgroup $\Gamma\le Q$, and a group homomorphism
$O:Q\to R$.
We call $O(q)$ the observation of the change $q$ and set
\begin{equation*}
K=\ker O,\qquad L=K\cap\Gamma
\tag{7.11}\label{eq:7.11}
\end{equation*}
Here $K$ consists of all the changes that do not appear in the observation, and $L$ of
those among them that are also compatible with the comparison.
The observation here is defined as a map preserving the composition of changes.

\end{definition}

\begin{theorem}[Determination by the observation alone]\label{thm:7.9}

The following conditions are equivalent.

\begin{enumerate}
\item There is a predicate $D:R\to\{\text{true},\text{false}\}$ such that
$q\in\Gamma\Longleftrightarrow D(O(q))$ for every $q\in Q$.
\item $K\subseteq\Gamma$.
\end{enumerate}

Moreover,
\begin{equation*}
O^{-1}(O(\Gamma))=\Gamma K
\tag{7.12}\label{eq:7.12}
\end{equation*}
always holds. Surjectivity of $O$ is not assumed.

\end{theorem}

\begin{proof}

Assume 1. Then $k\in K$ has the same observation as the identity element.
The identity element belongs to $\Gamma$, so $k\in\Gamma$, which gives 2.
Conversely, assume 2 and define $D(r)$ to be $r\in O(\Gamma)$.
If $O(q)=O(\gamma)$ with $\gamma\in\Gamma$, then
$k=\gamma^{-1}q\in K$ and $q=\gamma k\in\Gamma$.
This decomposition shows both inclusions of \eqref{eq:7.12} without assuming 2.

\end{proof}

What we have obtained here is the condition for compatibility to depend only on the observed
value.
To actually build a program that determines it, a way of computing $D$ is needed as well.

\begin{proposition}[Differences remaining within the same observation]\label{prop:7.10}

For $\gamma\in\Gamma$, the fiber of the observation and its compatible part are
\begin{equation*}
O^{-1}(O(\gamma))=\gamma K,\qquad
O^{-1}(O(\gamma))\cap\Gamma=\gamma L
\tag{7.13}\label{eq:7.13}
\end{equation*}
Specifying the base point $L$ in the set of left cosets $K/L=\{kL\mid k\in K\}$, we have
\begin{equation*}
\gamma k\in\Gamma\quad\Longleftrightarrow\quad kL=L.
\tag{7.14}\label{eq:7.14}
\end{equation*}
Hence $K/L$ being a single point is equivalent to the condition for the determination in
Theorem~\ref{thm:7.9}.

\end{proposition}

\begin{proof}

Agreement of the observations is equivalent to $\gamma^{-1}q\in K$.
Moreover, since $\gamma\in\Gamma$, the condition $\gamma k\in\Gamma$ is equivalent to
$k\in\Gamma$.
This gives \eqref{eq:7.13} and \eqref{eq:7.14}.
The set of cosets being a single point is equivalent to $K=L$, that is, to
$K\subseteq\Gamma$.

\end{proof}

$L$ is a normal subgroup of $\Gamma$, but need not be a normal subgroup of $K$.
What is needed here is therefore the set of cosets with a base point.

\begin{corollary}[The case of reading the coefficients only]\label{cor:7.11}

Fix an isomorphism comparison $c:X\xrightarrow{\sim}Y$ of full geometries.
From the groups fixing the base at the two endpoints, we take the homomorphisms reading the
coefficient component to be
$o_X:G_X^0\to\mathrm{Aut}(k_X)$ and $o_Y:G_Y^0\to\mathrm{Aut}(k_Y)$.
Write the coefficient isomorphism of the comparison as $\kappa:k_X\xrightarrow{\sim}k_Y$.
The conjugation $T(u)=cuc^{-1}$ satisfies
$o_Y(T(u))=\kappa o_X(u)\kappa^{-1}$.
Setting $K_X=\ker o_X$, $K_Y=\ker o_Y$, and $O=o_X\times o_Y$, we have
\begin{equation*}
\begin{gathered}
K=K_X\times K_Y,\qquad L=\{(a,T(a))\mid a\in K_X\},\\
K/L\xrightarrow{\sim}K_Y,\qquad [(a,b)]\longmapsto bT(a)^{-1}.
\end{gathered}
\tag{7.15}\label{eq:7.15}
\end{equation*}
The last correspondence is a bijection of sets sending the base point to the identity
element.
A necessary and sufficient condition for compatibility with the comparison to be
determinable from the observation of the coefficients alone is $K_Y=\{1\}$.

\end{corollary}

\begin{proof}

By the composition rule for the coefficient maps and Corollary~\ref{cor:7.5}, $T$ gives
$K_X\cong K_Y$, and $L$ is its graph.
Even if the representative of a coset is changed to $(a,b)(x,T(x))$, we have
$bT(x)T(ax)^{-1}=bT(a)^{-1}$, so \eqref{eq:7.15} is well defined.
The inverse is $y\mapsto[(1,y)]$.
Indeed, $(a,b)(a^{-1},T(a)^{-1})=(1,bT(a)^{-1})$, and both composites are the identity.
The condition for the determination follows from Proposition~\ref{prop:7.10}.

\end{proof}

Cases in which this set of cosets cannot be treated as a quotient group also occur in concrete cases.
If $K_X=K_Y=S_3$ (the group of permutations of three points) and $T=\mathrm{id}$, then $L$
is the diagonal subgroup.
For $s=(12)$ and $t=(23)$, the element
$(s,1)(t,t)(s,1)^{-1}=(sts^{-1},t)$ does not belong to the diagonal.
The bijection of \eqref{eq:7.15} is still valid in this case.

When the presentation of the input is carried as in Proposition~\ref{prop:7.7}, the
conjugations of the endpoint groups and of the coefficient groups commute with the
observation homomorphism.
Hence $K,L$ and the pointed cosets correspond as well.
Changing the presentation and regenerating does not change the conclusion as to whether
compatibility can be determined from the observation.

\section{Restriction to the normalization, reflection, and lifting}\label{sec:7.4}

When normalization is used as an observation, we study three properties separately.
That a change compatible in the original is compatible in the image as well is called
\textbf{preservation}, and that compatibility in the image implies compatibility in the
original is called \textbf{reflection}. Constructing an original of a compatible change
given in the image is \textbf{lifting}.

\begin{proposition}[Preservation of a comparison by a functor]\label{prop:7.12}

A functor $F:E\to D$ and a comparison $c:X\to Y$ induce the homomorphism
\begin{equation*}
\begin{gathered}
r_F:\mathrm{Aut}(X)\times\mathrm{Aut}(Y)
\longrightarrow\mathrm{Aut}(FX)\times\mathrm{Aut}(FY),\\
(u,v)\longmapsto(Fu,Fv)
\end{gathered}
\tag{7.16}\label{eq:7.16}
\end{equation*}
and we have $r_F(\Gamma_c)\subseteq\Gamma_{F(c)}$.
Applying this to the core normalization functor $\mathcal N$ of Chapter~\ref{chap:6} gives
preservation of an arbitrary comparison.
Moreover, the equality $\pi_N\mathcal N=\pi$ of the projections to the base in
Theorem~\ref{thm:6.21} yields a homomorphism between the comparison-preserving groups that
fix the base as well.
By Proposition~\ref{prop:6.22}, the same holds for the normalization functor of full
geometries.

\end{proposition}

\begin{proof}

Functors preserve isomorphisms, identities, and composition, so \eqref{eq:7.16} is a
homomorphism.
Applying $F$ to $vc=cu$ gives preservation.
Applying the equality on the base to $u,v$ shows that the condition of fixing the base is
preserved as well.

\end{proof}

Separately, a restriction can also be formed when idempotents are specified in an arbitrary
category.
In this case we choose the changes commuting with the normalization as the domain.

\begin{construction}[Restriction to the idempotent image]\label{cons:7.13}

Assume $dc=ce$ for a morphism $c:X\to Y$ and idempotents $e:X\to X$ and $d:Y\to Y$.
We take the centralizer groups and the comparison of the images to be
\begin{equation*}
\begin{aligned}
C_e&=\{u\in\mathrm{Aut}(X)\mid ue=eu\},\quad
C_d=\{v\in\mathrm{Aut}(Y)\mid vd=dv\},\\
H&=C_e\times C_d,\qquad
a=dce:(X,e)\longrightarrow(Y,d)
\end{aligned}
\tag{7.17}\label{eq:7.17}
\end{equation*}
Here $a$ is a morphism of $\mathrm{Kar}(E)$. The restriction is given by
\begin{equation*}
\begin{gathered}
r:H\longrightarrow\mathrm{Aut}_{\mathrm{Kar}(E)}(X,e)
\times\mathrm{Aut}_{\mathrm{Kar}(E)}(Y,d),\\
r(u,v)=(eue,dvd)
\end{gathered}
\tag{7.18}\label{eq:7.18}
\end{equation*}
The inverse of $eue$ is $eu^{-1}e$, and both composites are the identity morphism $e$ of
$(X,e)$.
By the commutation condition, the restriction preserves products as well.
Setting $\Gamma_0=H\cap\Gamma_c$ and $\Delta=\Gamma_a$, we have $r(\Gamma_0)\subseteq\Delta$.
Indeed, it suffices to sandwich $vc=cu$ between the idempotents and use $ue=eu$ and $vd=dv$.

Proposition~\ref{prop:7.12} treats all the automorphisms on which the normalization functor
is defined, and Construction~\ref{cons:7.13} treats the centralizer groups.
For the normalization functor of Chapter~\ref{chap:6}, the two homomorphisms agree on the centralizing changes.

In Theorem~\ref{thm:7.9} membership in the set of observed compatible changes could be used
for the determination.
Here the compatible subgroup on the side of the image is determined independently.
Besides the condition on the kernel, it is therefore also necessary that the compatible
changes on the side of the image that are actually reached be obtained from changes
compatible in the original. The second condition of the next theorem expresses this.

\end{construction}

\begin{theorem}[General determination of reflection and lifting]\label{thm:7.14}

Assume $r(\Gamma_0)\subseteq\Delta$ for a group homomorphism $r:H\to R$ and subgroups
$\Gamma_0\le H$ and $\Delta\le R$. Then
\begin{equation*}
r^{-1}(\Delta)=\Gamma_0
\quad\Longleftrightarrow\quad
\begin{cases}
\ker r\subseteq\Gamma_0,\\
r(\Gamma_0)=\Delta\cap\mathrm{im}\,r.
\end{cases}
\tag{7.19}\label{eq:7.19}
\end{equation*}
We take the restricted homomorphism to be $\bar r:\Gamma_0\to\Delta$ and its kernel to be
$L=\ker\bar r$.
A compatible lift of $\delta\in\Delta$ exists if and only if $\delta\in r(\Gamma_0)$.
Taking one of them to be $\gamma_0$, all the compatible lifts are
\begin{equation*}
\bar r^{-1}(\delta)=\gamma_0L
\tag{7.20}\label{eq:7.20}
\end{equation*}
and form a torsor under $L$ by multiplication on the right.
If $\bar r$ is surjective, the inclusion $L\hookrightarrow\Gamma_0$ gives the short exact
sequence
\begin{equation*}
1\longrightarrow L\longrightarrow\Gamma_0
\xrightarrow{\bar r}\Delta\longrightarrow1
\tag{7.21}\label{eq:7.21}
\end{equation*}
Here short exact means that the first map is injective, its image is the kernel of the next
map, and the last map is surjective.
When there is a group homomorphism $\sigma:\Delta\to\Gamma_0$ satisfying
$\bar r\sigma=\mathrm{id}_\Delta$, the sequence is said to split and $\sigma$ is called a
section.

\end{theorem}

\begin{proof}

If reflection holds, the preimage of the identity element lies in $\Gamma_0$, and every
image of $r$ belonging to $\Delta$ comes from $\Gamma_0$.
Conversely, assume the right-hand side and let $r(q)\in\Delta$.
By the condition on the image there is $\gamma\in\Gamma_0$ with $r(\gamma)=r(q)$, and
$\gamma^{-1}q\in\ker r\subseteq\Gamma_0$ gives $q\in\Gamma_0$.

If the lifts $\gamma_0,\gamma$ have the same observation, then $\gamma_0^{-1}\gamma\in L$.
This gives \eqref{eq:7.20}. The right action satisfies
$(\gamma\ell_1)\ell_2=\gamma(\ell_1\ell_2)$, and the element between two lifts is uniquely
determined as $\gamma_0^{-1}\gamma$.
Short exactness is just the definition of the kernel and surjectivity.

\end{proof}

It is $L=\ker\bar r$ that expresses the freedom of lifting.
What obstructs the reflection of compatibility, on the other hand, are those elements of the
kernel $\ker r$ of all endpoint changes that do not lie in $\Gamma_0$.
The two can be distinguished by the following examples.

\begin{example}[Lifts exist even when reflection fails]\label{ex:7.15}

In the category of finite sets we set $X=Y=\{0,1,2\}$ and $c=1_X$, and take $e=d$ to be the
constant map $x\mapsto0$.
The permutation $\tau=(12)$ commutes with $e$, but the endpoint pair $(\tau,1)$ does not
preserve the identity comparison.
Since the image is a one-element set, after the restriction both endpoints become the
identity, which preserves the comparison of the images.
The comparison-preserving group of the image is the trivial group, so its unique element
lifts to $(1,1)$.
Moreover, $(\tau,\tau)$ is a compatible lift of the same element as well.

\end{example}

\begin{example}[A case where a change of the image does not lift]\label{ex:7.16}

Using the same three-element set and $c=1_X$, we now take
\begin{equation*}
e(0)=0,\qquad e(1)=1,\qquad e(2)=1,
\qquad d=e
\tag{7.22}\label{eq:7.22}
\end{equation*}
Exchanging 0 and 1 at both endpoints of the image $\{0,1\}$ preserves the identity
comparison of the images.
A permutation $u$ commuting with $e$, however, puts $e^{-1}(0)$ and $e^{-1}(u(0))$ in
bijection.
The former has one element and $e^{-1}(1)$ has two, so $u(0)=1$ is impossible.
Since $e^{-1}(2)$ is empty, $u(0)=2$ is impossible as well.
Hence this exchange in the image has no lift inside the centralizer group.

These examples show that preservation, reflection, and the existence of lifts each require
their own proof.
In the next section we construct lifts using the concrete form of object formation and
normalization in AAT.

\end{example}

\section{Constructing lifts for full geometries}\label{sec:7.5}

A general idempotent can have fibers of different sizes, as in Example~\ref{ex:7.16}.
In the canonical normalization of AAT, an object splits into a configuration and the
structure data accompanying it.
Using this common form, a change given on the normalized image can be extended consistently
to all the original objects.

In this section we assume that the core $P$ of a full geometry $G$ satisfies
$\mathrm{Ad}(P)$ of Definition~\ref{def:5.39}.
That is, the normalization leaves the residuals of the equations and the coordinates of the
invariants and of the signature unchanged, and the identification of the types of operations
preserves the configuration maps.
We write $\mathcal A$ for the full subcategory of full geometries satisfying this assumption.
We use the normalization functor of Chapter~\ref{chap:6} as
$\mathcal N:\mathcal A\to\mathrm{Kar}(\mathcal A)$,
$\mathcal N(G)=(G,N_G)$, $\mathcal N(f)=fN_G$.

\subsection*{Carrying the data accompanying an object}

\begin{construction}[Extension preserving the selected objects]\label{cons:7.17}

We split an architecture object $A$ of Chapter~\ref{chap:1} into a configuration $C_A$ and
accompanying data $(S_A,Q_A,s_A,q_A)$. We write the set of all accompanying data, in a
universe of the required size, as
\begin{equation*}
D=\coprod_{S,Q\in\mathbb{U}}(S\times Q)
\tag{7.23}\label{eq:7.23}
\end{equation*}
Since the types $S,Q$ themselves are included in the data, all the objects are in one-to-one
correspondence with $\mathcal C\times D$.
Here $\mathcal C$ is the set of configurations.
We write the object selected by the core as $s_P(C)=(C,d_C)$ and fix a reference element
$d_*$ of $D$.
Such a $d_*$ is obtained using two one-element sets and their elements.

Let $\theta_C$ be the permutation exchanging $d_C$ and $d_*$.
When the two are equal we use the identity. Then
\begin{equation*}
\theta_C^2=1,\qquad \theta_C(d_C)=d_*,\qquad
\theta_C(d_*)=d_C.
\tag{7.24}\label{eq:7.24}
\end{equation*}
Let $\sigma$ be the Atom component of an automorphism $a$ of the normalized image, and
abbreviate the transport of a configuration as $\sigma C$.
If we simply take $(C,d)\mapsto(\sigma C,d)$, the selected object $(C,d_C)$ goes to
$(\sigma C,d_C)$, which need not agree with the selected object $(\sigma C,d_{\sigma C})$ at
the target.
We therefore combine the two exchanges, at the start and at the target, and define the
extension to all the objects by
\begin{equation*}
\widehat a(C,d)=
\bigl(\sigma C,\theta_{\sigma C}\theta_C(d)\bigr)
\tag{7.25}\label{eq:7.25}
\end{equation*}
This first aligns the selected value at the start with the common reference, and returns to
the selected value at the target.
Hence $\widehat a(s_P(C))=s_P(\sigma C)$, so object formation is preserved.
The same formula can be defined for a self-map of Atoms as well.

With this construction, the intermediate permutations cancel when changes are performed in
succession.
For the two extensions corresponding to the Atom components $\sigma,\rho$, we have
\begin{equation*}
\theta_{\rho\sigma C}\theta_{\sigma C}
\theta_{\sigma C}\theta_C
=\theta_{\rho\sigma C}\theta_C.
\tag{7.26}\label{eq:7.26}
\end{equation*}
Hence the map on objects preserves identities and composition.
That this equality holds not only for the selected objects but for all the accompanying data
is the reason we obtain a lift as a group homomorphism.

\end{construction}

\subsection*{Extending the components of the operations and of the geometry}

\begin{lemma}[A section for the automorphisms of a full geometry]\label{lem:7.18}

For every $G\in\mathcal A$ there exists a group homomorphism
\begin{equation*}
\sigma_G:\mathrm{Aut}(\mathcal N G)\longrightarrow\mathrm{Aut}(G),
\qquad \mathcal N(\sigma_G(a))=a
\tag{7.27}\label{eq:7.27}
\end{equation*}
This section keeps the base and the coefficient components, and its values commute with
$N_G$.

\end{lemma}

\begin{proof}

We represent a morphism of the normalized image by a morphism $f:G\to G$ of the original
category with $N_GfN_G=f$.
For objects we use Construction~\ref{cons:7.17}. We abbreviate the normalization of objects
of Chapter~\ref{chap:6} as $n=n_P$ and write the identification of the types of operations as
$\nu_{A,B}:\mathrm{Op}(A,B)\xrightarrow{\sim}\mathrm{Op}(nA,nB)$.
These identifications, which follow the equalities of the types, are the identity on the
selected objects and are compatible with the composition of successive identifications.
Object formation gives $f(nA)=n\widehat f(A)$, so we define the action on operations by
\begin{equation*}
S_G(f)_{A,B}^{\mathrm{op}}
=\nu_{\widehat f(A),\widehat f(B)}^{-1}
\,f_{nA,nB}^{\mathrm{op}}\,\nu_{A,B}
\tag{7.28}\label{eq:7.28}
\end{equation*}
For the identification of the intermediate objects we use the equality of object formation
above.
This is the map that normalizes at the start, applies the given morphism, and returns to the
type at the target of the extension.
Since $\nu$ and $f$ each preserve the action on configurations, so does this composite.

For the maps of the Atoms, the contexts, the equation indices, the invariant indices, and
the axes and their value ranges we use the components of $f$.
Preservation of the residuals of the equations follows by replacing $A$ with $nA$ at the
start, applying the preservation law of $f$, and returning $n\widehat f(A)$ to
$\widehat f(A)$ at the target.
The same three-step computation is carried out for the coordinates of the invariants and of
the signature.
The base and the coefficients, the covers and the overlaps, and the local maps of the
supports, the axes, and the observables are taken over from those of $f$.
Since the maps of Atoms and contexts to which these refer are unchanged, the equalities of
the raw system and the commutativity of the local restrictions are preserved as well.
Hence $S_G(f)$ is a morphism of full geometries.

By \eqref{eq:7.26} for objects and the cancellation of the intermediate $\nu\nu^{-1}$ in
\eqref{eq:7.28} for operations, we obtain, for endomorphisms $f,g$ of the normalized image,
\begin{equation*}
S_G(N_G)=1_G,\qquad
S_G(gf)=S_G(g)S_G(f),\qquad
S_G(f)N_G=N_GS_G(f)=f
\tag{7.29}\label{eq:7.29}
\end{equation*}
In the last equality we also use that $f=N_GfN_G$ gives $fN_G=N_Gf=f$.
The other components are taken over from $f,g$, so they satisfy the same equalities.
In particular, applying $S_G$ to the morphisms in both directions of an automorphism of
the normalized image, we find that both composites are $1_G$.
Taking this to be $\sigma_G$, we obtain \eqref{eq:7.27} and the homomorphism property.
Keeping the base and the coefficients follows from the construction, and commutation with
$N_G$ follows from \eqref{eq:7.29}.

\end{proof}

\subsection*{Combining the two endpoints so as to preserve the comparison}

Sections built separately at the two endpoints of a comparison need not be compatible with
that comparison.
We therefore lift only one side and determine the other by conjugation with the original
comparison.

\begin{theorem}[Compatible lifts for an isomorphism comparison]\label{thm:7.19}

Take $X,Y\in\mathcal A$ and an isomorphism $c:X\xrightarrow{\sim}Y$.
The restriction
$\bar r_N:\Gamma_c\to\Gamma_{\mathcal N(c)}$ of Proposition~\ref{prop:7.12} has the section
\begin{equation*}
\sigma_c(a,b)=\bigl(\sigma_X(a),\ c\sigma_X(a)c^{-1}\bigr)
\tag{7.30}\label{eq:7.30}
\end{equation*}
The base and coefficient components at the two endpoints agree with those of the given
$a,b$.
Hence the restriction to the subgroups fixing the base has a section as well.

\end{theorem}

\begin{proof}

Since the second component is built by conjugation, \eqref{eq:7.30} preserves the original
comparison $c$.
The normalization of the first component is $a$.
By $(a,b)\in\Gamma_{\mathcal N(c)}$ and functoriality, the normalization of the second
component is
\begin{equation*}
\mathcal N(c)\,a\,\mathcal N(c)^{-1}=b
\tag{7.31}\label{eq:7.31}
\end{equation*}
This gives $\bar r_N\sigma_c=\mathrm{id}$.
Both the section on the first component and the conjugation are group homomorphisms, so
$\sigma_c$ is a group homomorphism.
Since the normalization does not change the base and coefficient components, keeping them
for the second component follows from \eqref{eq:7.31}.

\end{proof}

\section{Information remaining in a generated comparison, and all the lifts}\label{sec:7.6}

Given a section, a compatible change in the original can be built from a compatible change
after normalization.
Not every change with the same value after normalization is compatible, however.
This difference can be confirmed by building an actual automorphism on all the objects of
AAT.

\begin{construction}[A nonidentity change that does not appear in the normalization]\label{cons:7.20}

Let $G\in\mathcal A$ and use the presentation $(C,d)$ of Construction~\ref{cons:7.17}.
The set $D$ has at least three distinct elements.
For instance, it suffices to fix $S=\{0,1,2\}$ and $Q=\{*\}$ and use the structure values
$s=0,1,2$.
For each $C$, choose two elements different from the selected value $d_C$ and let $\rho_C$
be the permutation exchanging them.
The map on objects
\begin{equation*}
\tau_G(C,d)=(C,\rho_C(d))
\tag{7.32}\label{eq:7.32}
\end{equation*}
fixes the configurations and all the selected objects, and applying it twice gives the
identity.
For instance, a configuration specified by the core has two objects that are exchanged, so
this map is not the identity.

This map can be extended to an automorphism of the full geometry.
For the operations we use $\nu_{\tau_G A,\tau_G B}^{-1}\nu_{A,B}$, and for the maps of the
Atoms, the base, the contexts, the equation indices, the invariant indices, the axes, the
coefficients, and the local geometry we use the identity.
The normalizations of the two objects agree, so the preservation conditions follow by the
same computation through the normalization as in Lemma~\ref{lem:7.18}.
The exchange of the accompanying data and the identification of the types of operations
return after two applications. Hence
\begin{equation*}
\begin{gathered}
\tau_G\ne1_G,\qquad \tau_G^2=1_G,\\
N_G\tau_G=\tau_GN_G=N_G,\qquad
\mathcal N(\tau_G)=1_{\mathcal N G}.
\end{gathered}
\tag{7.33}\label{eq:7.33}
\end{equation*}
This change fixes the base and the coefficients.
While it moves the distinction between objects, it does not appear in the components read
through the normalization under Ad.

\end{construction}

\begin{theorem}[Distinguishing and lifting for a generated comparison]\label{thm:7.21}

Fix the common input of Construction~\ref{cons:5.37}. From the square realized by finite
codes, the core and the strong edges at each vertex, the conditions on the end of each face
and on the two paths in the base, the specified comparisons, the chosen face, the
coefficients, and the original full geometry, we generate the canonical comparison
$\alpha:X\xrightarrow{\sim}Y$ of full geometries.
We assume Ad for the source core and use the canonical normalizations at the two endpoints.
By Lemma~\ref{lem:6.15} and Theorem~\ref{thm:6.16} of Chapter~\ref{chap:6}, the two
endpoints satisfy Ad as well, and $N_Y\alpha=\alpha N_X$.

We set $\Gamma=\Gamma_\alpha$ and $\Delta=\Gamma_{\mathcal N(\alpha)}$, take
$r_N$ to be the homomorphism on the full endpoint groups, and set $L=\ker(r_N|_\Gamma)$. Then:

\begin{enumerate}
\item $r_N|_\Gamma$ has a section $\sigma_\alpha$, and we obtain the split short exact
sequence \eqref{eq:7.34} below.
\item For every $\delta\in\Delta$, all the compatible lifts are
$\sigma_\alpha(\delta)L$, a right $L$-torsor.
\item The full endpoint groups also contain incompatible changes with the same value
$\delta$ after normalization.
Hence compatibility with $\alpha$ cannot be determined from the value after normalization
alone.
\end{enumerate}
\begin{equation*}
1\longrightarrow L\longrightarrow\Gamma
\xrightarrow{\bar r_N}\Delta\longrightarrow1
\tag{7.34}\label{eq:7.34}
\end{equation*}
These also hold for the groups fixing the base at the two endpoints.

\end{theorem}

\begin{proof}

1 and 2 follow from Theorems~\ref{thm:7.19} and 7.14.
For 3, set $q_-=(\tau_X,1_Y)$. By Construction~\ref{cons:7.20}, $r_N(q_-)=1$.
If $q_-\in\Gamma$, then $\alpha=\alpha\tau_X$, and invertibility of $\alpha$ gives the
contradiction $\tau_X=1_X$.
Hence
\begin{equation*}
q_+=\sigma_\alpha(\delta),\qquad
q_{\mathrm{bad}}=\sigma_\alpha(\delta)q_-
\tag{7.35}\label{eq:7.35}
\end{equation*}
have the same value $\delta$, and only the former is compatible.
If the latter were compatible as well, their difference $q_-$ would lie in the subgroup
$\Gamma$.
The section keeps the base component and $q_-$ fixes the base, so the same proof applies to
the base-fixing version.

\end{proof}

In the case of reading the coefficients only, Construction~\ref{cons:7.20} also gives an
actual incompatible pair.
The pair $(\tau_X,\alpha\tau_X\alpha^{-1})$ is compatible and $(\tau_X,1_Y)$ is not, yet the
observations of the base and the coefficients of both pairs are the same identity pair.
Hence the loss of information in Corollary~\ref{cor:7.11} occurs for comparisons of full
geometries generated from the same input as well.

\begin{corollary}[Selection of the normalization by the diagnosis]\label{cor:7.22}

Take the common input of Construction~\ref{cons:5.37} and a cochain on the same diagram, and
write the selection condition of Construction~\ref{cons:5.40} as $\chi$.
In the notation of Chapter~\ref{chap:6} we set $\beta=d\alpha=\alpha e$ and let $H$ be the
group of endpoint changes commuting with $e,d$.
The restriction of Construction~\ref{cons:7.13}, from $\Gamma_\alpha\cap H$ to
the comparison-preserving group $\Gamma_{\beta:(X,e)\to(Y,d)}$ of the image, has a
section.
Preservation and reflection of compatibility with $\alpha$ under this restriction are as
follows.

\begin{xltabular}{\linewidth}{@{}LLLL@{}}
\toprule
Normalization & Preservation of the comparison & Reflection of compatibility with $\alpha$ &
Lifting of the compatible changes of the image \\
\midrule
\endhead
$\chi$ is false, $e=1_X,d=1_Y$ & Holds & Holds & A section exists \\
\addlinespace[3pt]
$\chi$ is true, $e=N_X,d=N_Y$ & Holds & Fails & A section exists \\
\addlinespace[3pt]
The canonical normalization is chosen directly under Ad & Holds & Fails & A section exists \\
\bottomrule
\end{xltabular}

Each row holds for the subgroups fixing the base as well.

\end{corollary}

\begin{proof}

If the selection condition is false, the restriction is to the identity idempotents, so the
group is the same as the original comparison group.
In the true case, Theorem~\ref{thm:6.16} gives $e=N_X,d=N_Y$.
The element $\sigma_X(a)$ of Lemma~\ref{lem:7.18} commutes with $N_X$.
From $\alpha N_X=N_Y\alpha$, the conjugate $\alpha\sigma_X(a)\alpha^{-1}$ commutes with
$N_Y$ as well.
Hence \eqref{eq:7.30} gives a section inside $H$.
The failure of reflection follows from $(\tau_X,1_Y)$, which belongs to $H$ as well.
The proof is the same when the canonical normalization is chosen directly. The condition on
the base is kept in every construction.

\end{proof}

What is not reflected here is compatibility with $\alpha$ before normalization.
On the centralizer group $H$, preserving $\beta$ in the image is equivalent to preserving
$\beta$ as a morphism of the original category, and
\begin{equation*}
r^{-1}\bigl(\Gamma_{\beta:(X,e)\to(Y,d)}\bigr)
=H\cap\Gamma_{\beta:X\to Y}
\tag{7.36}\label{eq:7.36}
\end{equation*}
Indeed, by $d\beta=\beta=\beta e$ and the centralizing condition, the equality
$(dvd)\beta=\beta(eue)$ in the image reduces to $v\beta=\beta u$.
The normalization moves to a comparison that has forgotten part of the condition for
preserving the original comparison.

\section{Classification of changes by the connections of operations}\label{sec:7.7}

We bring the classification so far back to the operations of software.
An operation that carries part of the state unchanged ties together the changes that are
invisible to the readout.
The update of a lens, and a protocol that proceeds while keeping an internal state, have
this same form.

\begin{definition}[Systems of operations keeping hidden states]\label{def:7.23}

Take a directed multigraph $Q=(V,E,s,t)$, a set $K$, and a group
$H\le\mathrm{Aut}(Q)$ of the permitted visible changes.
An automorphism $u$ of the graph consists of a permutation $u_V$ of the vertices and a
permutation $u_E$ of the named edges, and we take it to preserve starts and ends.
We define the states, the observation, and the operation of each edge by
\begin{equation*}
S=V\times K,\qquad o(v,k)=v,\qquad
T_e(s(e),k)=(t(e),k)
\tag{7.37}\label{eq:7.37}
\end{equation*}
Here $K$ is the internal state that the operations do not change.
We take a state change following a visible change $u\in H$ to be a bijection
$h:S\xrightarrow{\sim}S$ satisfying $oh=u_Vo$.
The condition of preserving the named operations is, moreover, $hT_e=T_{u_E(e)}h$ on the
fiber over the start of each edge.
The directions and the names of the edges are kept as the input of this equality.

We write $B_Q$ for the group of all pairs $(u,h)$ following the readout, and $A_Q$ for the
subgroup of those that also preserve the named operations.
The observation of a change is the visible projection $(u,h)\mapsto u$.
These correspond to all the changes, the compatible changes, and the observation of
Theorem~\ref{thm:7.9}.
In what follows we study between which vertices the condition of preserving the operations
propagates, and determine this subgroup concretely.

\end{definition}

\begin{theorem}[Preservation of operations and connected components]\label{thm:7.24}

A state change following a fixed $u$ can be written uniquely, using permutations for each
vertex ($\mathrm{Sym}(K)$ is the group of all permutations of $K$), as
\begin{equation*}
h(v,k)=(u_V(v),\phi_v(k)),\qquad \phi_v\in\mathrm{Sym}(K)
\tag{7.38}\label{eq:7.38}
\end{equation*}
A necessary and sufficient condition for preserving the operations is
\begin{equation*}
\phi_{s(e)}=\phi_{t(e)}\qquad(e\in E)
\tag{7.39}\label{eq:7.39}
\end{equation*}
Hence the changes preserving the operations are in one-to-one correspondence with the
choices of one permutation for each element of the set $\pi_0(Q)$ of connected components
with the directions forgotten.

\end{theorem}

\begin{proof}

By the condition on the observation, $h$ sends $\{v\}\times K$ to $\{u_V(v)\}\times K$.
The inverse map also returns to the corresponding fiber, so the second component is a
permutation $\phi_v$.
Evaluating both sides of the preservation of operations at $(s(e),k)$, the first components
are $u_V(t(e))$ in either case, and the second components are $\phi_{t(e)}(k)$ and
$\phi_{s(e)}(k)$ respectively.
This gives \eqref{eq:7.39}.

The quotient of $V$ by the equivalence relation generated by the relation $s(e)\sim t(e)$
between vertices is $\pi_0(Q)$.
Since \eqref{eq:7.39} is preserved under reflexivity, symmetry, and transitivity, the family
of permutations descends uniquely to this quotient.
Conversely, a family of permutations indexed by the connected components satisfies
\eqref{eq:7.39} once pulled back to the vertices.

\end{proof}

Since the operation of an edge carries the internal state as it is, the same permutation is
required at its two ends.
The connected components express the range over which this equality propagates regardless of
the directions of the edges.

\begin{theorem}[Split short exact sequence and the freedom of changes]\label{thm:7.25}

We take the projection from the group $A_Q$ of changes preserving the operations to the
visible changes to be $p:A_Q\to H$.
The composition is $(u,h)(u',h')=(uu',hh')$, and for the families of \eqref{eq:7.38} it is
\begin{equation*}
(\phi\star\phi')_v=\phi_{u'_V(v)}\phi'_v
\tag{7.40}\label{eq:7.40}
\end{equation*}
Setting $M_Q=\prod_{j\in\pi_0(Q)}\mathrm{Sym}(K)$, we obtain the split short exact sequence
\begin{equation*}
1\longrightarrow M_Q\longrightarrow A_Q
\xrightarrow{p}H\longrightarrow1
\tag{7.41}\label{eq:7.41}
\end{equation*}
The section is $\sigma(u)=(u,h_u)$ with $h_u(v,k)=(u_V(v),k)$.
The set of changes over each $u\in H$ is a right torsor under $M_Q$.
For finite $V,K$, the number of such changes is
\begin{equation*}
|p^{-1}(u)|=(|K|!)^{|\pi_0(Q)|}
\tag{7.42}\label{eq:7.42}
\end{equation*}
\end{theorem}

\begin{proof}

The equalities for the observation and the operations are preserved under composition and
inverses, so $A_Q$ is a group.
By Theorem~\ref{thm:7.24}, the fiber over $u=1$ is isomorphic as a group to $M_Q$.
Since $\sigma(u)$ preserves the operations and satisfies $\sigma(uu')=\sigma(u)\sigma(u')$
and $p\sigma=\mathrm{id}$, the short exact sequence splits.
The right torsor structure of each fiber is obtained by applying Theorem~\ref{thm:7.14} to
$p$.
Counting the families of permutations of Theorem~\ref{thm:7.24} gives \eqref{eq:7.42}.

\end{proof}

\begin{corollary}[The version preserving the selected states]\label{cor:7.26}

Fix $k_0\in K$ and take $s_0(v)=(v,k_0)$ as the selected state.
If $hs_0=s_0u_V$ is required as well, the permutation of each connected component is
restricted to $\mathrm{Sym}(K)_{k_0}=\{\phi\mid\phi(k_0)=k_0\}$.
We obtain the bijection between this family of permutations and each fiber, the split short
exact sequence, and the right torsor under the kernel.
The section is the same $\sigma(u)$, which does not change the internal state.
In the finite case, the number of elements of each fiber is $((|K|-1)!)^{|\pi_0(Q)|}$.

\end{corollary}

\begin{proof}

Evaluating \eqref{eq:7.38} at $(v,k_0)$, the additional condition is $\phi_v(k_0)=k_0$.
The identity action on the hidden state satisfies this condition, so the same section can be
used.

\end{proof}

Using this classification, we find the condition under which preservation of the operations
can be determined from the visible change alone.
The kernel of the visible projection of the group $B_Q$ of changes following the readout is
$\prod_{v\in V}\mathrm{Sym}(K)$, and the part of it preserving the operations is $M_Q$.
Hence, by Theorem~\ref{thm:7.9}, preservation of the operations can be determined from the
visible change alone exactly when these two agree.
If there is an edge joining distinct vertices and $K$ has two or more elements, choosing
different permutations at its two ends gives a counterexample.
Reading the operations brings out the constraint forgotten by the visible projection.

\section{Consequences for lenses and protocols}\label{sec:7.8}

In Chapter~\ref{chap:1} we defined the states, the operations, and the
semantics-preserving morphisms of lenses and protocols independently, and put them in
correspondence with the construction of AAT having the roles of objects and typed operations.
Here we use that construction for the one-to-one correspondence of states linking the
situations before and after a refactoring.

Write the states before and after the change as $C,C'$, the reads as $g,g'$, and the
operations as $T_e,T'_e$.
Once the view and the names of the operations are fixed, what we require of a state
correspondence $h:C\xrightarrow{\sim}C'$ is $g'h=g$ and $hT_e=T'_eh$.
The former expresses the agreement of the results of the reads, and the latter the agreement
of ``convert and then operate'' with ``operate and then convert.''

Once one compatible correspondence $h_0$ has been obtained, every other correspondence can
be written uniquely as $h=h_0a$ using $a=h_0^{-1}h\in\mathrm{Sym}(C)$.
The condition on the reads then becomes $ga=g$ and the condition on the operations becomes
$aT_e=T_ea$.
Indeed, it suffices to compose $h_0^{-1}$ on the left of $h_0aT_e=T'_eh_0a=h_0T_ea$.
Once the states before and after are expressed in the same form through a reference
correspondence in this way, the automorphism groups of the previous section classify the
options for the state correspondence that remain in the refactoring.

\subsection*{Lenses: the total update ties all the views together}

\begin{proposition}[Changes of a product lens]\label{prop:7.27}

Fix a reference view $v_0\in V$ and a finite set $K$, and consider the self-changes of the
same product lens $C=V\times K$.
Fix a visible permutation $u:V\xrightarrow{\sim}V$.
The preservation conditions for a bijection $h:C\xrightarrow{\sim}C$ are
\begin{equation*}
gh=ug,\qquad h(p(c,w))=p(h(c),u(w))
\tag{7.43}\label{eq:7.43}
\end{equation*}
In the graph whose vertices are the views and whose named edges of updates are the ordered
pairs $(v,w)$, these changes correspond bijectively to the state changes that follow the
visible permutation $u$ and preserve all the named operations.
Each change has a unique permutation $\phi\in\mathrm{Sym}(K)$ with
\begin{equation*}
h(v,k)=(u(v),\phi(k))
\tag{7.44}\label{eq:7.44}
\end{equation*}
Conversely, this formula builds a change from an arbitrary $\phi$.
Hence there are $|K|!$ changes following a fixed visible permutation.

\end{proposition}

\begin{proof}

The read of a product lens is the first component, and the named edge $(v,w)$ sends
$(v,k)$ to $(w,k)$. Reading the two preservation conditions in their respective types gives
the conditions for a following state change.
Since the updates for all the ordered pairs are preserved, the hidden permutation at each
view agrees with the permutation at the reference view.
This is \eqref{eq:7.44}, and evaluating at the reference view gives uniqueness as well.
Conversely, $u\times\phi$ satisfies both preservation equations.
From these two inverses we obtain the bijection between the set of changes and
$\mathrm{Sym}(K)$, and the number.

\end{proof}

When the selected section is fixed as $s_0(v)=(v,k_0)$, the condition $hs_0=s_0u$ is
equivalent to the permutation of each fiber fixing $k_0$.

\begin{example}[Carrying over the payment data even when the editing of the shipping
address is separated]\label{ex:7.28}

In the refactoring of the opening, we separate the processing that edits the shipping
address from the keeping of the payment data.
The \texttt{Order(address, payment)} before the change is moved to the pair
\texttt{OrderDetails(address)} and \texttt{PaymentDetails(payment)} after the change. This
repackaging is the reference state correspondence.
We take the shipping addresses to be $V=\{0,1\}$ and the internal code of the payment data,
which is not displayed, to be $K=\{0,1,2\}$.
Views 0 and 1 are the two shipping addresses, home and work, and $k_0=0$ represents the
unset value of the payment data.
The correspondence between the old and the new codes is also part of the proposed change,
and the correspondence of the unset value 0 is fixed.
This is the product lens whose read returns the shipping address and whose update rewrites
only the shipping address and carries over the payment data.

We fix the visible change to be the identity and align the states before and after into the
same product decomposition.
The correspondences preserving the read of the shipping address and the unset value are the
four given by choosing at each view whether to exchange 1 and 2.
The correspondences that also preserve the update are the two that use the same
correspondence of codes at both shipping addresses.

For instance, suppose the conversion branches on the shipping address, keeping the code at
view 0 while exchanging 1 and 2 only at view 1. For this correspondence $h$,
\begin{equation*}
h(p((0,1),1))=(1,2),\qquad
p(h((0,1)),1)=(1,1)
\tag{7.45}\label{eq:7.45}
\end{equation*}
Changing the shipping address first and then moving to the new format gives code 2, while
moving to the new format and then changing the shipping address gives code 1, so the payment
data carried over does not match.
A test of the display of the shipping address and a test on an order whose payment data is
unset both miss this difference.
What is needed is to check the update and the conversion of states in both orders for an
order that already has payment data.

\end{example}

\subsection*{Protocols: independent changes remain on independent components}

\begin{proposition}[Protocols keeping an internal state]\label{prop:7.29}

Take the finite protocol graph $Q$ of Definition~\ref{def:1.36} and a finite set $K$, and
take the state at each vertex to be $F(v)=K$ and the action of each generating edge to be
$F(e)=\mathrm{id}_K$.
We take the target of the observation to be the constant one-element set. All parallel paths
then give the same identity map, so we obtain a realization satisfying the specified path
equations $\Pi$.

Take the group $H\le\mathrm{Aut}(Q)$ of visible changes preserving the path equivalence
generated by $\Pi$.
The invertible changes of the realization following a fixed $u\in H$ are the vertexwise
bijections $\phi_v:K\to K$ satisfying
\begin{equation*}
\phi_{t(e)}F(e)=F(u_E(e))\phi_{s(e)}
\tag{7.46}\label{eq:7.46}
\end{equation*}
and are the same as the operation-preserving changes of Definition~\ref{def:7.23}.
The classification of Theorem~\ref{thm:7.25} by the split short exact sequence, the section,
and the torsor under the kernel therefore applies to the group of changes.

\end{proposition}

\begin{proof}

We gather all the states of the realization, keeping the control points distinct, into
$\coprod_{v\in V}F(v)=V\times K$.
The readout of the control point is the observation $o(v,k)=v$ of the previous section.
Equation~\eqref{eq:7.46} is equivalent to $\phi_{t(e)}=\phi_{s(e)}$.
By the same induction on paths as in Proposition~\ref{prop:1.38}, this equality extends to
all the executions.
The visible changes preserve the path equivalence as well, so the correspondence is
determined on the quotient category too.
Conversely, an invertible change of the realization returns to this family of permutations
once the vertex components are read.
Both directions preserve composition, so the groups and the fibers over each visible change
correspond.

\end{proof}

The changes with $u=1$ over a fixed graph are the automorphisms of this realization in
$\mathrm{Prot}(Q,\Pi,1)$, where the observation is constant to a one-element set.
When $u$ is moved as well, we record at the same time the isomorphisms following the
reindexing of the control points and of the names of the operations.
For protocols with general state transitions, the actual edge maps remain inside
\eqref{eq:7.46}.
The classification above by the connected components alone is a consequence obtained from
the input that each edge carries the internal state by the identity.

\begin{example}[Gathering the execution contexts of workers]\label{ex:7.30}

Suppose two workers A and B with the same role each receive and execute jobs independently.
Consider a refactoring that gathers the arguments scattered inside each worker into a data
type \texttt{JobContext} and tidies up the handoff from the queue to the execution.
Even when a job proceeds to Running, the context in which it runs must not change.

We represent by the internal code $K=\{0,1\}$ one component of the context whose handoff we examine.
A permutation is the choice of how to match the old and the new representations of the same
context.
Among the four control points $V=\{0,1,2,3\}$, we take 0 and 1 to be Queued and Running for
A, and 2 and 3 to be Queued and Running for B.
The edges $a:0\to1$ and $b:2\to3$ representing the start of execution carry this code as it
is.
In this model there is no edge handing a context between the workers, so it has two
connected components.

\end{example}

\begin{figure}[htbp]
\centering
% Figure 7.1 -- context-preserving refactoring in two workers.
% Same content as the Japanese manuscript's figure ../../figures/ch07-refactoring-workers.svg, redrawn in TikZ.
\begin{tikzpicture}[
    x=1mm, y=1mm,
    state/.style={circle, draw, semithick, minimum size=8mm, inner sep=0pt, fill=white},
    box/.style={rounded corners=2mm, draw=black!45, semithick, fill=black!4},
    arr/.style={-{Stealth[length=5pt]}, semithick},
    lbl/.style={font=\small},
    note/.style={font=\small}
  ]
  % Worker A
  \draw[box] (0,0) rectangle (62,36);
  \node[font=\small\bfseries] at (31,32) {Worker A};
  \node[state] (a0) at (12,20) {$0$};
  \node[state] (a1) at (50,20) {$1$};
  \draw[arr] (a0) -- (a1) node[midway, above, lbl] {$a$: start job}
                          node[midway, below, lbl] {preserves $k$};
  \node[lbl] at (12,11) {Queued};
  \node[lbl] at (50,11) {Running};
  \node[lbl] at (12,6)  {$\varphi_0$};
  \node[lbl] at (50,6)  {$\varphi_1$};
  \node[lbl] at (31,2.5) {Required: $\varphi_0 = \varphi_1$};
  % Worker B
  \begin{scope}[xshift=68mm]
    \draw[box] (0,0) rectangle (62,36);
    \node[font=\small\bfseries] at (31,32) {Worker B};
    \node[state] (b2) at (12,20) {$2$};
    \node[state] (b3) at (50,20) {$3$};
    \draw[arr] (b2) -- (b3) node[midway, above, lbl] {$b$: start job}
                            node[midway, below, lbl] {preserves $k$};
    \node[lbl] at (12,11) {Queued};
    \node[lbl] at (50,11) {Running};
    \node[lbl] at (12,6)  {$\varphi_2$};
    \node[lbl] at (50,6)  {$\varphi_3$};
    \node[lbl] at (31,2.5) {Required: $\varphi_2 = \varphi_3$};
  \end{scope}
  \node[note] at (65,-7)
    {$K = \{0,1\}$: identity or swap per worker $\to$ $2 \times 2 = 4$ choices};
\end{tikzpicture}

\caption{\textbf{Aligning the correspondence at both ends of a handoff.} If the visible
change is the identity, $\phi_0=\phi_1$ is required for A and $\phi_2=\phi_3$ for B.
If, for instance, the codes are exchanged on the queue side only while the execution side
uses the previous correspondence, then it can still be confirmed that a job proceeds from
Queued to Running, but the correspondence of the contexts breaks.
Since there is no handoff between the two workers, the internal representation of A and that
of B can be chosen independently.}\label{fig:7.1}
\end{figure}

We compare the case where the placement of the workers is fixed with the case where A and B,
which have the same role, are exchanged.

\begin{xltabular}{\linewidth}{@{}LLL@{}}
\toprule
Condition checked & Placement of the workers fixed & Workers A and B exchanged \\
\midrule
\endhead
Following the readout of the control points & $2^4=16$ & $2^4=16$ \\
\addlinespace[3pt]
Named operations preserved as well & $2^2=4$ & $2^2=4$ \\
\bottomrule
\end{xltabular}

The exchange moves the vertices $0\leftrightarrow2$ and $1\leftrightarrow3$ and the edges
$a\leftrightarrow b$ at the same time.
One exchange that does not change the internal state can be chosen.
Composing it with the permutations of the two components gives the four
operation-preserving changes following the exchange.

\subsection*{Connection with the comparison-preserving group}

For lenses, restricting the correspondence of typed morphisms of
Proposition~\ref{prop:1.43} to invertible morphisms identifies the group of
semantics-preserving automorphisms for a fixed view with the group of typed automorphisms on
the side of AAT.
For protocols as well, the same proposition makes the invertible morphisms preserving the
vertex maps, the named operations, and the observation correspond.
Fix an adapter $c:X\to Y$ and apply Theorem~\ref{thm:4.41} to the fully faithful functor of
Proposition~\ref{prop:1.43}.
The group of endpoint pairs satisfying $vc=cu$ is then identified between lenses over a
fixed view, or protocols over a fixed schema, and their typed constructions.
This correspondence is compatible with the projections to the endpoints.

The equalities for the operations determine which subgroup of the comparison-preserving
group is chosen, and restrict the freedom of lifting.
For lenses the total update ties all the views together, and in the example of protocols a
change remains for each independent worker.
For product lenses we describe this freedom by the bijection and the number over a fixed
visible permutation, and for protocols by the split short exact sequence and each fiber.

\section*{Summary of the Chapter}

In this chapter we distinguished the determination, the construction, and the classification
of all the candidates for changes preserving a comparison, and found the information needed
for each.

\begin{itemize}
\item \textbf{The group of changes preserving a comparison.} For a single comparison $c$, we
constructed the group of endpoint pairs satisfying $vc=cu$.
For the base-fixing groups of full geometries, its projections and kernels give the
condition for following and the freedom among all the candidates.
Changes generated from a common original object preserve the comparison and are compatible
with the re-presentation of the input by isomorphisms
(\S\S\ref{sec:7.1}--\ref{sec:7.2}).
\item \textbf{Determination by observation.} A necessary and sufficient condition for
compatibility to be determinable from the observation $O$ of changes alone is
$\ker O\subseteq\Gamma$.
When a change that does not appear in the observation breaks the comparison, a compatible
change and an incompatible change have the same observation (\S\ref{sec:7.3}).
\item \textbf{Construction of lifts and reflection.} For the canonical normalization of AAT,
a section can be constructed from the common presentation of the accompanying data.
For the same generated comparison, even when compatible lifts exist, the compatibility of a
given change cannot in some cases be determined from the value after normalization alone
(\S\S\ref{sec:7.4}--\ref{sec:7.5}).
\item \textbf{The freedom among all the lifts.} A section gives one compatible lift, and the
kernel of the restricted homomorphism describes the other candidates.
Each fiber is a torsor under this kernel, and the loss of information and the options for
lifting can be read from the same restriction map (\S\ref{sec:7.6}).
\item \textbf{Classification by operations.} We described the state changes preserving the
operations for each connected component, and obtained a split short exact sequence onto the
visible group.
For product lenses we put the changes following a single visible permutation in
correspondence with the hidden permutations.
For protocols keeping an internal state, independent changes remain on independent
components (\S\S\ref{sec:7.7}--\ref{sec:7.8}).
\end{itemize}

In the refactoring of the opening, we require that the payment data be carried over by the
same correspondence after the shipping address has been updated.
In the example of a lens, the four candidates preserving the display of the shipping address
and the unset value are narrowed down to the two that also preserve updates.
In the example of workers as well, the sixteen candidates preserving the correspondence of
the control points are narrowed down to the four that preserve the handoff of contexts.
These numbers show how much freedom of change remains once each operation ties the
correspondences between states together.

In Chapter~\ref{chap:8} we move on to the problem of assembling the original structure and
the structure-preserving morphisms from the local data read off from objects and changes.
The classification of changes in this chapter gives concretely the morphisms that this
reconstruction must keep and the freedom remaining in the observation.
